\documentclass[12pt,a4paper]{ctexart}

\usepackage{geometry}
\geometry{left=2.5cm,right=2.5cm,top=2.5cm,bottom=2.5cm}
\usepackage{hyperref}
\usepackage{booktabs}
\usepackage{xcolor}
\usepackage{enumitem}
\usepackage{listings}
\usepackage{fancyhdr}

\pagestyle{fancy}
\fancyhf{}
\fancyhead[L]{宝宝日历 · 智能育儿看板}
\fancyhead[R]{项目文档}
\fancyfoot[C]{\thepage}
\renewcommand{\headrulewidth}{0.4pt}

\title{\textbf{宝宝日历 · 智能育儿看板}\\[4pt]\large 项目文档 v2.1}
\author{dede2024-cmd}
\date{\today}

\begin{document}

\maketitle
\thispagestyle{empty}
\newpage

\tableofcontents
\newpage

\section{项目概述}

宝宝日历（Baby Calendar）是一个专为婴幼儿家庭设计的智能育儿看板，以 HTML 单页应用形式运行，无需安装任何 App。它在一块屏幕（iPad / 手机 / 平板 / 电脑桌面）上集中展示\textbf{作息时间表、早教主题、天气信息}三大模块，并支持夜间自动切换极简模式。

项目完全免费开源，部署在 GitHub Pages 上，任何能打开浏览器的设备均可使用。

\subsection{核心用户}

\begin{itemize}
    \item 带孩子的\textbf{阿姨、保姆、祖辈}——抬眼即知当前该做什么
    \item 孩子\textbf{父母}——远程了解家中日程安排
    \item 家庭中的\textbf{多设备场景}——不同房间的屏幕各取所需
\end{itemize}

\subsection{项目地址}

\begin{itemize}
    \item 主仓库：\url{https://github.com/dede2024-cmd/baby-calendar}
    \item 桌面版仓库：\url{https://github.com/dede2024-cmd/baby-calendar-desktop}
    \item 在线访问：\url{https://dede2024-cmd.github.io/baby-calendar/}
\end{itemize}

\section{功能特性}

\subsection{作息时间表}

按照「一苇、二杭」两个宝宝的日常节奏，内置11个时段作息项，涵盖喂奶、用餐、户外活动、午睡、早教、洗澡、入睡等。每个时段带有类型标签：

\begin{itemize}
    \item \textbf{必须}（红色）——喂奶、用餐、睡眠等固定事项
    \item \textbf{灵活}（黄色）——户外活动等可调整事项
    \item \textbf{早教}（紫色）——主题早教时段
    \item \textbf{建议}（绿色）——洗澡等日常护理
\end{itemize}

\textbf{当前时间段自动高亮}，视觉上清晰区分「现在该做什么」。

\subsection{今日早教主题}

按星期轮换7大主题：

\begin{table}[htbp]
\centering
\begin{tabular}{cll}
\toprule
\textbf{星期} & \textbf{主题} & \textbf{能力培养} \\
\midrule
日 & 亲子时光 & 情感联结、安全感 \\
一 & 颜色认知 & 视觉辨识、分类能力 \\
二 & 动物世界 & 认知记忆、语言模仿 \\
三 & 身体认知 & 自我意识、语言表达 \\
四 & 音乐律动 & 节奏感知、大运动协调 \\
五 & 自然探索 & 感官体验、观察力 \\
六 & 生活技能 & 自理能力、秩序感 \\
\bottomrule
\end{tabular}
\caption{早教主题轮换表}
\end{table}

每个主题配有4个具体活动卡片，分别对应早晨、户外、主题早教时段、睡前四个场景。

\subsection{实时天气}

接入 Open-Meteo 免费天气 API，展示合肥市庐阳区的天气信息：

\begin{itemize}
    \item 当前室外气温（每10分钟自动刷新）
    \item 当日最高温 / 最低温
    \item 天气状况（晴/多云/雨等，含对应 emoji 图标）
    \item 降水概率
\end{itemize}

\subsection{夜间模式（22:00—次日 6:00）}

晚上10点后自动切换至极简夜间模式：

\begin{itemize}
    \item 只显示\textbf{当前室外温度 + 时间}，居中超大字体
    \item 夜间主题为深黑色背景（\#020202）+ 暗灰色文字，不刺眼
    \item 作息表和早教内容全部隐藏
    \item 早上6点自动恢复完整显示
\end{itemize}

\subsection{多设备自适应}

通过 CSS 媒体查询实现单个网页适配所有设备：

\begin{table}[htbp]
\centering
\begin{tabular}{p{2.5cm}p{3cm}p{4cm}}
\toprule
\textbf{设备类型} & \textbf{分辨率} & \textbf{触发布局} \\
\midrule
1K 平板（iPad 3） & 1024×768 & 横屏5栏顶栏 + 3栏内容区 \\
2K 平板（iPad Pro） & 2048×1536+ & 同上，字体等比适配 \\
手机（荣耀/华为） & ～390×844 & 竖屏窄版，垂直堆叠布局 \\
桌面浏览器 & 任意 & 横屏大屏，信息密度最高 \\
\bottomrule
\end{tabular}
\caption{设备适配表}
\end{table}

\subsection{iPad 2 专用版}

针对 iPad 2（iOS 9.3.5，Safari 版本老旧）提供独立兼容版本：

\begin{itemize}
    \item 全部使用传统 JavaScript（无 ES6+ 语法）和基础 CSS（无变量、无 flexbox）
    \item 浅色主题（\#fffdf7 米白背景）
    \item 固定 1024×768 布局
    \item 访问地址：\texttt{ipad2.html}
\end{itemize}

\subsection{桌面壁纸模式}

通过 Plash（macOS）或 Lively Wallpaper（Windows）可将日历设为桌面壁纸，开机自启、常驻显示。

\section{技术架构}

\subsection{技术栈}

\begin{table}[htbp]
\centering
\begin{tabular}{llp{5.5cm}}
\toprule
\textbf{层} & \textbf{技术} & \textbf{说明} \\
\midrule
前端框架 & 无（原生 HTML/CSS/JS） & 单文件，零依赖，加载极快 \\
天气数据 & Open-Meteo API & 免费、无需注册、全球覆盖 \\
部署 & GitHub Pages & 免费静态托管，自动 HTTPS \\
本地服务 & Python HTTP Server & 可选，局域网直连 \\
壁纸引擎 & Plash / Lively Wallpaper & 将网页设为操作系统壁纸 \\
\bottomrule
\end{tabular}
\caption{技术栈一览}
\end{table}

\subsection{代码结构}

\begin{lstlisting}[basicstyle=\ttfamily\small]
calendar-server/
|-- index.html          # 主页面（自适应深色主题，~730行）
|-- ipad2.html          # iPad 2 兼容版（浅色主题，~520行）
|-- index.html.bak.*    # 自动备份文件
\end{lstlisting}

主页面核心逻辑：

\begin{itemize}
    \item \textbf{天气获取}：\texttt{fetchWeather()} 通过 Open-Meteo API 异步拉取，每10分钟刷新
    \item \textbf{夜间检测}：\texttt{checkNightMode()} 每秒检查当前时间，切换 CSS 类
    \item \textbf{作息高亮}：\texttt{renderSchedule()} 根据当前时间匹配作息项
    \item \textbf{主题轮换}：\texttt{renderEdu()} 根据星期几选择对应早教主题
\end{itemize}

\subsection{浏览器兼容性}

\begin{table}[htbp]
\centering
\begin{tabular}{lcc}
\toprule
\textbf{浏览器} & \textbf{主版本} & \textbf{iPad 2 版} \\
\midrule
Safari 14+（iOS 14+） & \checkmark & — \\
Safari 9（iOS 9） & × & \checkmark \\
Chrome 80+ & \checkmark & \checkmark \\
Firefox 80+ & \checkmark & \checkmark \\
Edge 80+ & \checkmark & \checkmark \\
微信内置浏览器 & \checkmark & \checkmark \\
\bottomrule
\end{tabular}
\caption{浏览器兼容性矩阵}
\end{table}

\section{部署指南}

\subsection{GitHub Pages 部署（推荐）}

无需任何服务器，GitHub 免费托管，自动 HTTPS：

\begin{enumerate}
    \item 将 \texttt{index.html} 上传至 GitHub 仓库根目录
    \item 仓库 Settings → Pages → Source 选择 \texttt{main} 分支 → 保存
    \item 等待1—2分钟，访问 \texttt{https://用户名.github.io/仓库名/}
\end{enumerate}

\subsection{局域网部署}

在家里的 Mac 上启动本地服务：

\begin{lstlisting}[basicstyle=\ttfamily\small]
cd ~/calendar-server
python3 -m http.server 8080 --bind 0.0.0.0
\end{lstlisting}

访问地址：\texttt{http://[Mac的局域网IP]:8080/index.html}

\subsection{各设备访问方式}

\begin{itemize}
    \item \textbf{iPad / iPhone}：Safari 打开 → 底部「分享」→「添加到主屏幕」
    \item \textbf{安卓手机}：Chrome 打开 → 右上角 ⋮ →「添加到主屏幕」
    \item \textbf{macOS 桌面壁纸}：安装 Plash → Add Website → 输入桌面版地址
    \item \textbf{Windows 桌面壁纸}：安装 Lively Wallpaper → 添加网址
\end{itemize}

\section{定制指南}

\subsection{修改作息时间}

编辑 \texttt{index.html} 中的 \texttt{schedule} 数组：

\begin{lstlisting}[basicstyle=\ttfamily\small]
{ time: "08:40", task: "早餐：鸡蛋 + 粥", type: "red", range: [8,40] }
\end{lstlisting}

\texttt{type} 取值：\texttt{red}（必须）/\texttt{yellow}（灵活）/\texttt{green}（建议）/\texttt{purple}（早教）。

\subsection{修改早教主题}

编辑 \texttt{eduData} 数组，每天一个主题对象，含 \texttt{title}、\texttt{ability}、\texttt{activities}。

\subsection{修改天气城市}

修改 JS 中的经纬度常量：

\begin{lstlisting}[basicstyle=\ttfamily\small]
const LAT = 31.878, LON = 117.265, CITY = "合肥·庐阳";
\end{lstlisting}

\subsection{调整夜间模式时间}

在 \texttt{checkNightMode()} 函数中修改判断条件：

\begin{lstlisting}[basicstyle=\ttfamily\small]
var isNight = (h >= 22 || h < 6); // 22点至次日6点
\end{lstlisting}

\section{常见问题}

\begin{description}
    \item[Q：天气数据不显示？] A：检查网络是否连通。页面会显示默认温度（25°C）作为兜底。
    \item[Q：iPad 上地址栏去不掉？] A：必须通过「添加到主屏幕」方式打开才会全屏。
    \item[Q：老 iPad 显示白屏？] A：使用 \texttt{ipad2.html} 专用版，兼容 iOS 9。
    \item[Q：如何同时更新多个仓库？] A：\texttt{baby-calendar} 和 \texttt{baby-calendar-desktop} 需同步更新。
\end{description}

\section{版本历史}

\begin{table}[htbp]
\centering
\begin{tabular}{cll}
\toprule
\textbf{日期} & \textbf{版本} & \textbf{更新内容} \\
\midrule
2026-06-16 & v2.1 & 新增庐阳区天气、夜间模式、iPad 2 专用版 \\
2026-06-15 & v2.0 & 深色主题、多栏自适应布局、跑马灯 \\
2026-06-09 & v1.0 & 初始版本，浅色主题固定布局 \\
\bottomrule
\end{tabular}
\caption{版本历史}
\end{table}

\section{许可证}

本项目采用 MIT 许可证，可自由使用、修改和分发。

\end{document}
